\section{Number of Provinces}
\label{sec:graph-number-of-provinces}

\problemstatement{Given an $n \times n$ adjacency matrix
$\textit{isConnected}$ for $n$ cities (where
$\textit{isConnected}[i][j]=1$ means cities $i$ and $j$ are directly
connected), find the number of \textbf{provinces} -- groups of cities
connected directly or indirectly, where a city in one group has no
connection to a city in any other group.}

\begin{patternbox}
\emph{``Count the number of disconnected groups''} is exactly counting
\textbf{connected components}: run a traversal (DFS or BFS) from every
\textbf{unvisited} node, marking everything that traversal reaches.
Each time a traversal is launched from a fresh, unvisited node, that is
one new province discovered.
\end{patternbox}

\subsection*{Algorithm}

\begin{enumerate}
  \item Maintain a $\textit{visited}$ array and a $\textit{provinces}$
        counter, both initialized to $0$/false.
  \item For each city $i$ from $0$ to $n-1$: if $i$ is not visited,
        increment $\textit{provinces}$, then run a full DFS (or BFS)
        from $i$, marking every reachable city visited.
  \item Return $\textit{provinces}$.
\end{enumerate}

\begin{ideabox}[Why This Always Terminates with the Correct Count]
Every city belongs to exactly one province (provinces are, by
definition, a partition of the cities). The traversal from an
unvisited city visits \emph{exactly} that city's entire province (no
more, no less, since the adjacency matrix's connections define the
province's boundary). So each outer-loop iteration that finds an
unvisited city is necessarily the \emph{first} city of a province not
yet counted -- which is precisely why incrementing the counter there,
and nowhere else, gives the correct total.
\end{ideabox}

\subsection*{C++ implementation}

\begin{lstlisting}
#include <bits/stdc++.h>
using namespace std;

void dfs(int node, vector<vector<int>>& isConnected, vector<bool>& visited) {
    visited[node] = true;
    int n = isConnected.size();
    for (int neighbor = 0; neighbor < n; neighbor++) {
        if (isConnected[node][neighbor] == 1 && !visited[neighbor]) {
            dfs(neighbor, isConnected, visited);
        }
    }
}

int findCircleNum(vector<vector<int>>& isConnected) {
    int n = isConnected.size();
    vector<bool> visited(n, false);
    int provinces = 0;

    for (int i = 0; i < n; i++) {
        if (!visited[i]) {
            provinces++;
            dfs(i, isConnected, visited);
        }
    }
    return provinces;
}

int main() {
    vector<vector<int>> isConnected = {
        {1, 1, 0},
        {1, 1, 0},
        {0, 0, 1}
    };
    cout << findCircleNum(isConnected) << endl; // 2
    return 0;
}
\end{lstlisting}

\subsection*{Dry run}

\dryrun{Take the $3\times3$ matrix above (cities $0,1$ connected;
city $2$ isolated).}

$i=0$: unvisited, $\textit{provinces}=1$; DFS from $0$ visits $1$
(connected) and stops (city $2$ not connected to $0$ or $1$);
visited $=\{0,1\}$. $i=1$: already visited, skip. $i=2$: unvisited,
$\textit{provinces}=2$; DFS from $2$ visits only itself. Final answer:
$2$ provinces.

\begin{complexitybox}
\textbf{Time Complexity.} $O(n^2)$ -- the adjacency matrix
representation forces an $O(n)$ scan per node during DFS (checking
every possible neighbor), across $n$ nodes total.
\textbf{Space Complexity.} $O(n)$ for the \textit{visited} array and
recursion stack.
\end{complexitybox}

\begin{takeawaybox}
``Count the number of connected components'' is one of the most
frequently recurring graph templates: launch a traversal from every
unvisited node, counting launches. This exact loop structure --
\textit{not} the traversal itself -- is what changes across this
chapter's later problems (Number of Islands, Section~\ref{sec:graph-number-of-islands},
reuses this identical loop on an implicit grid graph instead of an
explicit matrix).
\end{takeawaybox}
